\documentclass{article}


\usepackage{tikz} 
\usepackage{tikz-cd}
\usetikzlibrary{automata, positioning, arrows} 

\usepackage{array}
\usepackage{amsthm}
\usepackage{amsfonts}
\usepackage{amsmath}
\usepackage{amssymb}
\usepackage{fullpage}
\usepackage{color}
\usepackage{parskip}
\usepackage{hyperref}
  \hypersetup{
    colorlinks = true,
    urlcolor = blue,       % color of external links using \href
    linkcolor= blue,       % color of internal links 
    citecolor= blue,       % color of links to bibliography
    filecolor= blue,        % color of file links
    }
    
\usepackage{listings}
\usepackage[utf8]{inputenc}                                                    
\usepackage[T1]{fontenc}                                                       

\definecolor{dkgreen}{rgb}{0,0.6,0}
\definecolor{gray}{rgb}{0.5,0.5,0.5}
\definecolor{mauve}{rgb}{0.58,0,0.82}

\lstset{frame=tb,
  language=haskell,
  aboveskip=3mm,
  belowskip=3mm,
  showstringspaces=false,
  columns=flexible,
  basicstyle={\small\ttfamily},
  numbers=none,
  numberstyle=\tiny\color{gray},
  keywordstyle=\color{blue},
  commentstyle=\color{dkgreen},
  stringstyle=\color{mauve},
  breaklines=true,
  breakatwhitespace=true,
  tabsize=3
}

\newtheoremstyle{theorem}
  {\topsep}   % ABOVESPACE
  {\topsep}   % BELOWSPACE
  {\itshape\/}  % BODYFONT
  {0pt}       % INDENT (empty value is the same as 0pt)
  {\bfseries} % HEADFONT
  {.}         % HEADPUNCT
  {5pt plus 1pt minus 1pt} % HEADSPACE
  {}          % CUSTOM-HEAD-SPEC
\theoremstyle{theorem} 
   \newtheorem{theorem}{Theorem}[section]
   \newtheorem{corollary}[theorem]{Corollary}
   \newtheorem{lemma}[theorem]{Lemma}
   \newtheorem{proposition}[theorem]{Proposition}
\theoremstyle{definition}
   \newtheorem{definition}[theorem]{Definition}
   \newtheorem{example}[theorem]{Example}
\theoremstyle{remark}    
  \newtheorem{remark}[theorem]{Remark}

\title{CPSC-354 Report}
\author{ Angelina House  \\ Chapman University}

\date{\today} 

\begin{document}

\maketitle

\begin{abstract}
This report documents my learning and progress throughout the Fall 2025 semester in CPSC-354: Programming Languages. 
\end{abstract}

\setcounter{tocdepth}{3}
\tableofcontents


\section{Introduction}\label{intro}

This report chronicles my work in CPSC-354: Programming Languages during Fall 2025. Each weekly subsection contains restated homework prompts, fully typeset solutions, and a brief reflection or question that ties the technical material to broader themes.



\section{Week by Week}\label{homework}

\subsection{Week 1}\label{sec:week1}
\smallskip

\subsection*{Homework}
Discuss the MU puzzle from \emph{Gödel, Escher, Bach} (Chapter 1; see course PDF pp.~41--42). Does the string \texttt{MU} have a derivation in the MIU formal system starting from \texttt{MI}? Why or why not?

\subsubsection*{The MU (MIU) System Recap}
We begin with the single axiom string \texttt{MI}. The system has four production rules:
\begin{enumerate}
  \item If a string ends in \texttt{I}, append a \texttt{U} at the end.
  \item If the string has the form \texttt{M}x, duplicate everything after \texttt{M}: \texttt{M}x $\to$ \texttt{M}xx.
  \item Replace any occurrence of \texttt{III} by \texttt{U}.
  \item Delete any occurrence of \texttt{UU}.
\end{enumerate}

From \texttt{MI}, we can derive many strings by applying the rules:
\[
\texttt{MI} \xRightarrow{(2)} \texttt{MII} \xRightarrow{(2)} \texttt{MIIII} \xRightarrow{(3)} \texttt{MUI}.
\]
This illustrates how the rules expand and compress substrings, but it remains unclear how to end with exactly \texttt{MU}.

Let $I(s)$ be the number of \texttt{I}'s in a string $s$. Initially, $I(\texttt{MI})=1$, i.e., $1\ (\mathrm{mod}\ 3)$. Examine how each rule affects $I$ modulo $3$:
\begin{itemize}
  \item Rule 1 adds a \texttt{U}; $I$ is unchanged.
  \item Rule 2 doubles the count of \texttt{I}'s: $I \mapsto 2I$. Thus if $I \equiv 1\pmod{3}$ then $2I \equiv 2\pmod{3}$.
  \item Rule 3 subtracts three \texttt{I}'s (\texttt{III} $\to$ \texttt{U}); the residue modulo $3$ is unchanged.
  \item Rule 4 deletes \texttt{U}'s only; $I$ is unchanged.
\end{itemize}
Therefore, starting from $I \equiv 1\pmod{3}$, any derivable string keeps $I \equiv 1$ or $I \equiv 2\pmod{3}$, but never $0\pmod{3}$. The target \texttt{MU} has $I(\texttt{MU})=0$ ($0\pmod{3}$), so it is \emph{not} derivable. Hence, the MU puzzle has no solution.

The puzzle shows how stepping outside a formal system and using an invariant (here, the count of \texttt{I}'s modulo 3) can decisively answer a reachability question that seems open-ended from within the rules.

\subsubsection*{Reading Summary \& Question}
\textbf{Summary.} Hofstadter introduces formal systems via the MIU puzzle and contrasts \emph{inside-the-system} rule manipulation with \emph{meta-level} reasoning using invariants. This frames later topics on meaning, proof, and limitations of formal systems.

\textbf{Question.} Are there any situations where the MU puzzle could be true, and how much would we have to change to get to that point?

\subsection{Week 2}\label{sec:week2}
\smallskip

\subsection*{Homework}
Using the course notes on abstract reduction systems (ARS), determine for each of the seven given ARSs whether it is \emph{terminating}, \emph{confluent}, and whether it \emph{has unique normal forms (UN)}. Then, fill the grid of the eight possible \((\text{Confluent}, \text{Terminating}, \text{Has UN})\) combinations with examples or explain why some combinations are impossible.

\subsubsection*{Quick Definitions}
\begin{itemize}
  \item \textbf{Terminating:} there is no infinite chain \(x_0 \to x_1 \to x_2 \to \cdots\).
  \item \textbf{Confluent:} if \(x \to^* y\) and \(x \to^* z\), then \(y\) and \(z\) are joinable.
  \item \textbf{Has unique normal forms (UN):} every element has a (and in fact the) unique normal form.
\end{itemize}

\subsubsection*{Seven ARSs: Diagrams \& Properties}
Below, each ARS is shown as a tiny diagram (compiled with \texttt{tikz-cd}) followed by its properties and a one–line justification.

\paragraph{1) \(A=\varnothing\).}
\[
\text{(empty set, no nodes or arrows)}
\]
\textbf{Props:} Terminating \(\checkmark\), Confluent \(\checkmark\), UN \(\checkmark\). \emph{Reason:} All properties hold vacuously.

\paragraph{2) \(A=\{a\},\ R=\varnothing\).}
\[
\begin{tikzcd}
a
\end{tikzcd}
\]
\textbf{Props:} Terminating \(\checkmark\), Confluent \(\checkmark\), UN \(\checkmark\). \emph{Reason:} No arrows; \(a\) is a (unique) normal form.

\paragraph{3) \(A=\{a\},\ R=\{(a,a)\}\).}
\[
\begin{tikzcd}
a \arrow[out=45,in=135,loop, ""]
\end{tikzcd}
\]
\textbf{Props:} Terminating \(\times\), Confluent \(\checkmark\), UN \(\times\). \emph{Reason:} Self–loop gives infinite chain; only one object so forks (if any) are trivially joinable; no normal form exists.

\paragraph{4) \(A=\{a,b,c\},\ R=\{(a,b),(a,c)\}\).}
\[
\begin{tikzcd}
& a \arrow[dl] \arrow[dr] & \\
b & & c
\end{tikzcd}
\]
\textbf{Props:} Terminating \(\checkmark\), Confluent \(\times\), UN \(\times\). \emph{Reason:} From \(a\) reach distinct NFs \(b\) and \(c\) that are not joinable.

\paragraph{5) \(A=\{a,b\},\ R=\{(a,a),(a,b)\}\).}
\[
\begin{tikzcd}
a \arrow[loop left] \arrow[r] & b
\end{tikzcd}
\]
\textbf{Props:} Terminating \(\times\), Confluent \(\checkmark\), UN \(\checkmark\). \emph{Reason:} There exists an infinite chain via the loop on \(a\), but all reductions from \(a\) can join at \(b\), and every element’s normal form is uniquely \(b\).

\paragraph{6) \(A=\{a,b,c\},\ R=\{(a,b),(b,b),(a,c)\}\).}
\[
\begin{tikzcd}
& a \arrow[dl] \arrow[dr] & \\
b \arrow[loop left] & & c
\end{tikzcd}
\]
\textbf{Props:} Terminating \(\times\), Confluent \(\times\), UN \(\times\). \emph{Reason:} Loop at \(b\) breaks termination; from \(a\) you can reach \(b\) (non–NF forever) or \(c\) (an NF), so confluence and UN fail.

\paragraph{7) \(A=\{a,b,c\},\ R=\{(a,b),(b,b),(a,c),(c,c)\}\).}
\[
\begin{tikzcd}
& a \arrow[dl] \arrow[dr] & \\
b \arrow[loop left] & & c \arrow[loop right]
\end{tikzcd}
\]
\textbf{Props:} Terminating \(\times\), Confluent \(\times\), UN \(\times\). \emph{Reason:} Loops at \(b\) and \(c\); from \(a\) diverge to non–joinable loops; no normal forms.

\subsubsection*{Summary Table (Seven ARSs)}
\[
\begin{array}{c|c|c|c|l}
\# & \text{Terminating} & \text{Confluent} & \text{Has UN} & \text{Quick reason} \\\hline
1 & \checkmark & \checkmark & \checkmark & \text{Empty set: vacuous} \\
2 & \checkmark & \checkmark & \checkmark & \text{Single node, no edges} \\
3 & \times & \checkmark & \times & \text{Self–loop; no NF} \\
4 & \checkmark & \times & \times & a\to b\ \text{or}\ a\to c\ \text{(two NFs)} \\
5 & \times & \checkmark & \checkmark & \text{Loop at }a\ \text{but unique NF }b \\
6 & \times & \times & \times & \text{Loop at }b;\ a\to b\ \text{vs }a\to c \\
7 & \times & \times & \times & \text{Loops at }b,c;\ a\to b\ \text{vs }a\to c \\
\end{array}
\]

\subsubsection*{Eight Combinations Grid}
\[
\begin{array}{c|c|c|l}
\text{Confluent} & \text{Terminating} & \text{Has UN} & \text{Example or reason} \\\hline
\text{T} & \text{T} & \text{T} & \text{Example \#2 (or \#1)} \\
\text{T} & \text{T} & \text{F} & \textbf{Impossible:}\ \text{Confluence+termination}\Rightarrow\text{unique NF} \\
\text{T} & \text{F} & \text{T} & \text{Example \#5} \\
\text{T} & \text{F} & \text{F} & \text{Example \#3} \\
\text{F} & \text{T} & \text{T} & \textbf{Impossible:}\ \text{UN in a terminating system}\Rightarrow\text{confluent} \\
\text{F} & \text{T} & \text{F} & \text{Example \#4} \\
\text{F} & \text{F} & \text{T} & \textbf{Impossible:}\ \text{UN (for all elements)}\Rightarrow\text{confluent} \\
\text{F} & \text{F} & \text{F} & \text{Example \#6 or \#7} \\
\end{array}
\]

\subsubsection*{Reflection \& Question}
\textbf{Question.} What is a real life example of something that has two possible end states like in ARS \#4 where $a$ can go to either $b$ or $c$?
%  end of semester work  

\section{Synthesis}


\section{Evidence of Participation}


\section{Conclusion}\label{conclusion}


\begin{thebibliography}{99}


\end{thebibliography}

\end{document}